\section{Transpiler soundness}
\label{sec:soundness}

The Z3 backend (\S\ref{sec:implementation}) is the link between
the operational semantics of \S\ref{sec:semantics} and the
empirical claims of \S\ref{sec:evaluation}: the contrast tool,
\texttt{narrow-defence}, and the constrained-contrast scorer of
\S\ref{subsec:case_law_diff} all interrogate Z3 models of the
encoded library and read back ``conviction'' truth values. A
reader has to be able to trust that the Z3 model and the
operational semantics agree --- that the scorer's headline
numbers measure something semantically grounded rather than an
artefact of the encoding. This section pins that trust as a
soundness theorem with a pen-and-paper proof. Mechanisation is
deferred to a follow-on artefact (\S\ref{sec:limitations}); the
working theorem is the structural-correspondence claim that
follows.

\subsection{Theorem statement}
\label{subsec:soundness-statement}

Let $\mathcal{Z}$ denote the Z3 generator's translation of a
module $M$: a finite set of Z3 boolean and integer constants
$\mathcal{V}_M$ together with a finite set of assertions
$\Phi_M$. For each statute $S \in M$ with section number $n$,
$\mathcal{Z}$ allocates two named booleans
$\mathit{el}_n \equiv \texttt{<n>\_elements\_satisfied}$ and
$\mathit{cv}_n \equiv \texttt{<n>\_conviction}$ in
$\mathcal{V}_M$, and asserts the biconditionals
$\mathit{el}_n \Leftrightarrow \bigwedge_i e_i$ and
$\mathit{cv}_n \Leftrightarrow (\mathit{el}_n \,\wedge\,
\bigwedge_\xi \neg \mathit{fires}_\xi)$ in $\Phi_M$, where the
$e_i$ are leaf-element booleans for $S$ and the
$\mathit{fires}_\xi$ are exception-fired booleans
(\S\ref{subsec:semantics-exceptions}).

A \emph{satisfying assignment} $\mathcal{M}$ for $\Phi_M$ assigns
each variable in $\mathcal{V}_M$ a value such that every
assertion in $\Phi_M$ evaluates to $\top$. Given $\mathcal{M}$,
let $F_\mathcal{M}$ denote the corresponding fact pattern: the
boolean truth value of every $e_i \in \mathcal{V}_M$ becomes the
fact-map entry $F_\mathcal{M}(x_i) = \mathcal{M}(e_i)$ for the
underlying element identifier $x_i$.

\begin{theorem}[Z3-Operational Soundness]
\label{thm:z3-sound}
For every module $M$, every statute $S = \langle n, \ldots
\rangle \in M$, and every satisfying assignment $\mathcal{M}$ of
$\Phi_M$:
\[
  \mathcal{M}(\mathit{cv}_n) = \top
  \;\;\Longleftrightarrow\;\;
  \langle F_\mathcal{M}, \Sigma_M \rangle
  \vdash S \Downarrow_{\!\textsc{c}} \top.
\]
\end{theorem}

In words: the Z3 conviction Bool is true in a model iff the
operational semantics of \S\ref{sec:semantics} convicts under the
fact pattern that model encodes. The contrast and
narrow-defence tools' headline numbers are therefore
measurements of operational-semantic agreement, not artefacts of
the SAT encoding.

\subsection{Proof structure}
\label{subsec:soundness-proof}

The proof decomposes into four lemmas, one per judgement of
\S\ref{sec:semantics}. Each lemma states a per-judgement
correspondence between the Z3 encoding and the operational
rules; the main theorem follows by composition.

\begin{lemma}[Element correspondence]
\label{lem:soundness-elem}
For every leaf element $e = \langle \kappa, x, \rho \rangle$ in
$S$ and satisfying assignment $\mathcal{M}$:
\[
  \mathcal{M}(e_x) = b
  \;\;\Longleftrightarrow\;\;
  \langle F_\mathcal{M}, \Sigma_M \rangle \vdash e
  \Downarrow b.
\]
\end{lemma}

\begin{proof}[Sketch]
Direct from the construction of $F_\mathcal{M}$: the fact-map
entry for $x$ is, by definition, $\mathcal{M}(e_x)$. Both
\textsc{Eval-Elem} (when $x \in \mathrm{dom}(F)$) and
\textsc{Eval-Elem-Missing} (when $x \notin \mathrm{dom}(F)$, in
which case the Z3 boolean is unconstrained but the constructed
$F_\mathcal{M}$ records its assigned truth value regardless)
agree with the looked-up value. The element-kind tag $\kappa$ is
inert at this layer.
\end{proof}

\begin{lemma}[Element-graph correspondence]
\label{lem:soundness-graph}
For every element-tree $\mathit{Es}$ in $S$ and satisfying
assignment $\mathcal{M}$:
\[
  \mathcal{M}(\mathit{el}_n) = b
  \;\;\Longleftrightarrow\;\;
  \langle F_\mathcal{M}, \Sigma_M \rangle \vdash
  \mathit{Es} \Downarrow_{\!\textsc{el}} b.
\]
\end{lemma}

\begin{proof}[Sketch]
Induction on the structure of $\mathit{Es}$. Base case: a leaf
element, by Lemma~\ref{lem:soundness-elem}. Inductive cases:

\textit{$\textsf{all\_of}$ group.} The Z3 generator emits
$\mathit{el}_n \Leftrightarrow \bigwedge_i b_i$ (or the
sub-group-level boolean's biconditional). The operational rule
\textsc{Eval-AllOf} computes the same conjunction. By the IH on
each member, both sides agree.

\textit{$\textsf{any\_of}$ group.} Symmetric, with $\bigvee$
replacing $\bigwedge$ on both sides.
\end{proof}

\begin{lemma}[Exception correspondence]
\label{lem:soundness-exc}
For every exception $\xi \in \mathit{Xs}$ and satisfying
assignment $\mathcal{M}$:
\[
  \mathcal{M}(\mathit{fires}_\xi) = b
  \;\;\Longleftrightarrow\;\;
  \langle F_\mathcal{M}, \Sigma_M \rangle \vdash \xi
  \Downarrow_{\!\textsc{f}} b.
\]
\end{lemma}

\begin{proof}[Sketch]
Induction on the priority DAG, well-founded by the
defeats-acyclicity invariant
(\S\ref{subsec:semantics-exceptions}: the linter rejects cyclic
\texttt{defeats} edges, so the DAG admits a topological sort;
induction proceeds along the reverse order).

\textit{Base case:} an exception with no in-degree in the
defeats relation. Its \textsc{Fires} premise simplifies to
``guard true'' alone, which the Z3 generator emits as the
single-conjunct biconditional $\mathit{fires}_\xi
\Leftrightarrow \mathcal{Z}(\gamma)$. By
Lemma~\ref{lem:soundness-graph} (extended to general boolean
expressions over the fact-map; the extension is mechanical), the
guard's Z3-truth and operational-truth coincide.

\textit{Inductive case:} an exception $\xi$ with non-empty
$\mathrm{defeats}^{-1}(\xi)$. The Z3 generator emits
$\mathit{fires}_\xi \Leftrightarrow \mathcal{Z}(\gamma_\xi)
\,\wedge\, \bigwedge_{\xi' \in \mathrm{defeats}^{-1}(\xi)}
\neg \mathit{fires}_{\xi'}$. The operational rule
\textsc{Fires} asserts the same conjunction (with
\textsc{NotFires-Defeated} closing the negative case). By the
IH on each $\xi'$ (which has strictly smaller priority-DAG depth)
and the guard correspondence, both sides agree.
\end{proof}

\begin{lemma}[Cross-section correspondence]
\label{lem:soundness-cross}
For every \texttt{is\_infringed}($n'$) or
\texttt{apply\_scope}($n'$, $F'$) reference appearing in any
guard or expression of any statute in $M$:
\[
  \mathcal{M}(\mathit{cv}_{n'}) = b
  \;\;\Longleftrightarrow\;\;
  \mathit{cross}(n', \cdot) \Downarrow b
\]
where $\mathit{cross}$ is the operational reduction
(\textsc{IsInfringed} or \textsc{ApplyScope},
\S\ref{subsec:semantics-cross-section}) under either the ambient
or supplied fact pattern.
\end{lemma}

\begin{proof}[Sketch]
The Z3 generator's \texttt{\_conviction\_bool($n'$)} helper
declares $\mathit{cv}_{n'}$ lazily and dedupes by name: any
reference to the same section number resolves to the identical
Z3 atom. Therefore, when statute $n'$ is itself translated by
$\mathcal{Z}$, the conviction-biconditional asserted for $n'$
constrains the same atom that the cross-section reference reads.
By Lemmas~\ref{lem:soundness-graph} and~\ref{lem:soundness-exc}
applied to the inner statute, $\mathit{cv}_{n'}$'s truth tracks
the inner conviction judgement. The
\textsc{IsInfringed}/\textsc{ApplyScope} rules of
\S\ref{subsec:semantics-cross-section} thread the same fact
pattern as that conviction judgement (ambient $F$ for
\texttt{is\_infringed}; supplied $F'$ for \texttt{apply\_scope}),
so the operational and Z3 readings coincide.

A subtle case: when the referenced section $n'$ is \emph{not}
in $M$ (cross-library reference), the Z3 atom $\mathit{cv}_{n'}$
exists but is unconstrained --- the SAT solver may freely
assign it. The operational rule on the same input is undefined
(the $\Sigma$ lookup fails). We treat the cross-library case as
out of scope for Theorem~\ref{thm:z3-sound}; the case-law
differential testing of \S\ref{subsec:case_law_diff}
deliberately uses only same-module fact patterns to stay within
the theorem's scope.
\end{proof}

\subsection{Main theorem}
\label{subsec:soundness-main}

\begin{proof}[of Theorem~\ref{thm:z3-sound}]
The Z3 generator's conviction biconditional for statute $S$
asserts
\[
  \mathit{cv}_n \Leftrightarrow
  \mathit{el}_n \,\wedge\,
  \bigwedge_\xi \neg \mathit{fires}_\xi.
\]

By Lemma~\ref{lem:soundness-graph}, $\mathcal{M}(\mathit{el}_n)$
agrees with the operational element judgement
$\mathit{Es} \Downarrow_{\!\textsc{el}} b_E$. By
Lemma~\ref{lem:soundness-exc}, each $\mathcal{M}(\mathit{fires}_\xi)$
agrees with the operational exception judgement $\xi
\Downarrow_{\!\textsc{f}} b_\xi$.
Lemma~\ref{lem:soundness-cross} extends the same correspondence
through any cross-section reference reachable from a guard.

Substituting into the biconditional:
\[
  \mathcal{M}(\mathit{cv}_n) =
  b_E \,\wedge\, \bigwedge_\xi \neg b_\xi.
\]

The operational rule \textsc{Conviction} computes precisely
this: $S \Downarrow_{\!\textsc{c}} \top$ iff the elements'
conjunction is $\top$ and no exception fires. The two sides are
therefore equal under any satisfying assignment, which is
Theorem~\ref{thm:z3-sound}.
\end{proof}

\subsection{Penalty correspondence and mechanisation status}
\label{subsec:soundness-penalty}

The penalty algebra of \S\ref{subsec:semantics-penalty} reduces
to range pairs $\langle \mathit{lo}, \mathit{hi} \rangle$. The Z3
generator emits range bounds as integer-arithmetic assertions
(\texttt{imp\_min} $\le$ \texttt{imp\_max} and so on). The
correspondence here is structural rather than predicate-level:
the rules of \S\ref{subsec:semantics-penalty} produce a
\emph{value} (a range), whereas the Z3 encoding emits
\emph{constraints} on integer variables. Both sides agree on the
satisfying-range membership test (an imprisonment of $d$ days
satisfies the constraints iff
$\mathit{lo} \le d \le \mathit{hi}$), and that suffices for the
empirical claims --- the contrast and narrow-defence tools
exercise conviction (Theorem~\ref{thm:z3-sound}), not penalty
ranges.

\paragraph{Mechanisation.} A fully mechanised proof in Coq
or Lean would replace each lemma sketch above with a kernel-
checked derivation. The natural target is the
\texttt{\_conviction\_bool} helper's correspondence as a
simulation argument between the AST evaluator and the Z3
generator's encoding function, both written as Coq functions over
the AST datatype. We estimate three to six months of effort for a
competent proof-assistant user starting from the pen-and-paper
proof above; Catala's comparable mechanisation of its default-
calculus correctness theorem is the proximate model. We treat
the mechanisation as an artefact-level deliverable left for a
follow-on effort, and \S\ref{sec:limitations} carries the
explicit ``unmechanised'' caveat alongside the rest of the
paper's known unverified claims.
